%% -----------------------------------------------------------------
%% 北京邮电大学研究生学位论文 - 中文摘要
%% 格式要求（2025年1月15日规范）：
%% 1. “摘要”标题黑体三号字，居中对齐，首行无缩进，段前0行，段后1行，单倍行距，中间空1个全角空格
%% 2. 摘要内容用四号宋体，两端对齐，首行缩进2字符，段前0行，段后0行，固定值20磅
%% 3. “关键词：”黑体四号字，其后关键词为宋体四号，两端对齐，首行缩进2字符，段前1行，段后0行，固定值20磅
%% 4. 关键词3～5个，每一关键词之间用全角分号隔开（；），最后一个关键词后不打标点符号
%% -----------------------------------------------------------------

\begin{cnabstract}
摘要应简明表达学位论文的内容要点，体现研究工作的核心思想。重点说明本项科研的目的和意义、研究方法、研究成果、结论，注意突出具有创新性的成果和新见解的部分。文字应精练、确切，中文摘要字数一般在千字左右。

本论文针对当前信息通信与智能计算领域的关键科学与工程问题，深入开展了理论建模、算法设计与系统验证研究。首先，系统梳理了国内外相关领域的研究现状与技术瓶颈；其次，提出了一种高鲁棒性、低时延的协同处理机制与自适应优化框架；接着，通过严谨的数学推导论证了所提方法在收敛速度与系统容量上的理论界限；最后，基于大规模真实数据集与原型平台开展了充分的对比实验，全面评估了系统在不同网络负载与动态环境下的运行表现。

实验结果表明，所提出的创新方案能够显著降低计算开销，提升吞吐量与决策精度，满足新一代信息处理系统对高可靠性与超低时延的严苛指标要求。本文的工作为该领域的进一步理论演进与工业应用提供了坚实的理论依据与参考方案。

\cnkeywords{北京邮电大学；研究生学位论文；\LaTeX{}模板；格式规范；信息与通信工程}
\end{cnabstract}
